\chapter{La Simulation de Crise (Exercice Tabletop)}
\label{chap:simcrise}

La théorie est utile, mais la gestion de crise s'apprend par la pratique. Ce chapitre est dédié à un exercice de simulation de type \textbf{Tabletop} (exercice sur table). Il ne s'agit pas de tests techniques informatiques, mais de mise en situation des processus de décision et de communication.

\section{Présentation de l'Exercice}
% ------------------------------------------------------------------------------

\subsection{Objectifs Pédagogiques}

\begin{itemize}
    \item Tester la montée en puissance de la Cellule de Crise.
    \item Exercer la prise de décision sous pression et avec une information incomplète.
    \item Pratiquer la coordination entre les cellules (Stratégique, Opérationnel, Communication).
    \item Identifier les failles du Plan de Gestion de Crise (PGC).
\end{itemize}

\subsection{Mécanique du Jeu (Time Compression)}

L'exercice se déroule sur 1h30 en classe, mais simule une journée de crise (24h).
Le formateur joue le rôle du "Maître du Jeu" (Game Master) et injecte des événements (Twists) qui changent la donne.

% ------------------------------------------------------------------------------
\section{Phase 0 : Briefing et Attribution des Rôles (15 min)}
% ------------------------------------------------------------------------------

\subsection{Distribution des Rôles}

Les étudiants sont répartis en sous-groupes ou jouent individuellement les rôles clés définis au Chapitre 2.

\begin{table}[htbp]
\centering
\begin{tabular}{p{4cm}p{8cm}}
    \toprule
    \textbf{Rôle (Étudiant)} & \textbf{Objectif Personnel} \\
    \midrule
    \textbf{PDG (Directeur de Crise)} & Préserver l'entreprise, trancher les arbitrages difficiles. \\
    \textbf{DSI (Resp. Opérationnel)} & Restaurer le SI, fournir des faits techniques fiables. \\
    \textbf{Dir. Com (Porte-parole)} & Protéger l'image, rédiger les communiqués. \\
    \textbf{Juriste (DPO)} & Assurer la conformité (RGPD), protéger les dirigeants pénalement. \\
    \textbf{DRH} & Gérer le stress interne, empêcher les fuites salariés. \\
    \textbf{Scribe} & Noter la timeline exacte. \\
    \bottomrule
\end{tabular}
\caption{Matrice des rôles pour le Tabletop}
\end{table}

\subsection{Contexte Initial}

Le Maître du Jeu lit l'introduction :
\begin{quote}
"Nous sommes vendredi 17h00. LiZbiZ s'apprête à fermer pour le week-end. Soudain, le téléphone du DSI sonne : tous les écrans de l'entrepôt affichent un message rouge 'Your files are encrypted'. Le téléphone du standard commence à sonner : un client tweet qu'il a reçu un email de rançon avec ses données personnelles en pièce jointe."
\end{quote}

% ------------------------------------------------------------------------------
\section{Phase 1 : La Réaction Immédiate (20 min)}
% ------------------------------------------------------------------------------

\textbf{Temps simulé :} Vendredi 17h15.

\subsection{Consignes pour les équipes}
\begin{itemize}
    \item \textbf{Stratégie :} Décidez-vous d'éteindre les serveurs ? Prévenez-vous la police ?
    \item \textbf{Communication :} Que dites-vous aux employés qui rentrent chez eux ?
    \item \textbf{Opérationnel :} Quels sont les premiers diagnostics ?
\end{itemize}

\subsection{Injection d'Événement (Twist 1)}
Le Maître du Jeu annonce :
\begin{quote}
"Un journaliste du 'Figaro' appelle le standard. Il a connaissance de la rançon et demande si LiZbiZ a l'intention de payer. La réponse est attendue dans 10 minutes."
\end{quote}

% ------------------------------------------------------------------------------
\section{Phase 2 : L'Escalade (20 min)}
% ------------------------------------------------------------------------------

\textbf{Temps simulé :} Vendredi 20h00.

\subsection{Injection d'Événement (Twist 2)}
\begin{quote}
"La CNIL appelle la direction. Ils ont été alertés par un collectif de défense des consommateurs. Ils demandent le rapport de violation complet pour lundi matin."
\end{quote}

\subsection{Consignes pour les équipes}
\begin{itemize}
    \item \textbf{Juridique :} Que répondez-vous à la CNIL ? Avez-vous toutes les infos requises (Article 33 RGPD) ?
    \item \textbf{Technique :} L'attaquant a-t-il accédé les données de santé (risque élevé) ou juste les noms ?
\end{itemize}

% ------------------------------------------------------------------------------
\section{Phase 3 : La Tempête Médiatique (20 min)}
% ------------------------------------------------------------------------------

\textbf{Temps simulé :} Samedi 10h00.

\subsection{Injection d'Événement (Twist 3)}
\begin{quote}
"Le hashtag \#LiZbiZDataLeak est en tendance sur Twitter (X). Une fausse rumeur circule : 'La banque de LiZbiZ a fait faillite, tout l'argent est bloqué'. Les clients paniquent et demandent des remboursements massifs."
\end{quote}

\subsection{Consignes pour les équipes}
\begin{itemize}
    \item \textbf{Communication :} Organisez une conférence de presse simulée (Questions/Réponses avec le reste de la classe jouant les journalistes).
    \item \textbf{Stratégie :} Démentez-vous la rumeur bancaire ? Comment séparer le problème cyber du problème financier ?
\end{itemize}

% ------------------------------------------------------------------------------
\section{Phase 4 : Résolution et Debriefing (15 min)}
% ------------------------------------------------------------------------------

\textbf{Temps simulé :} Lundi 09h00. L'ERP est restauré à partir des sauvegardes saines.

\subsection{Fin de l'Exercice}

Le Maître du Jeu arrête le chronomètre. L'exercice s'arrête.

\subsection{Le "Hotwash" (RETEX à chaud)}

C'est le moment le plus important de l'apprentissage. On passe en revue ce qui s'est passé.

\textbf{Questions de debriefing :}
\begin{enumerate}
    \item \textbf{Decision Making :} Les décisions ont-elles été prises rapidement ? Manquait-on d'informations ?
    \item \textbf{Communication :} Le message a-t-il été cohérent ? A-t-on dit des choses contradictoires ?
    \item \textbf{Rôles :} Le PDG a-t-il été débordé par la technique ? Le DSI a-t-il été trop long à diagnostiquer ?
    \item \textbf{Leçons :} Qu'est-ce qu'on ferait différemment la prochaine fois ?
\end{enumerate}

\begin{boxalert}
\textbf{Conclusion typique :} On réalise souvent que la Cellule de Crise a passé trop de temps à chercher des informations techniques précises (Qui ? Comment ?) et pas assez à gérer l'émotion des clients et des employés. La leçon : En crise, l'humain et la communication priment sur la technique pure.
\end{boxalert}

% ==============================================================================
% Fichier : 05_module_crise.tex (Extrait Chapitre 5)
% Description : Le Retour d'Expérience (RETEX)
% ==============================================================================